\documentclass[12pt]{article}

% Paquetes básicos
\usepackage[utf8]{inputenc}
\usepackage[T1]{fontenc}
\usepackage[spanish]{babel}
\usepackage{amsmath, amssymb}
\usepackage[margin=2.5cm]{geometry}
\usepackage{hyperref}
\usepackage{coffeestains}
\usepackage{tikzducks}
\title{Propuesta de Simposio}
\author{Sara Sofía Rocha Ortiz \\ Universidad del Valle}
\date{\today}

\begin{document}
	\maketitle
	
	\section*{Tema }
	
	Reproducción númerica de la dinámica de un pantinador con restricciones no holónomas, mostrando la integrabilidad y transición al comportamiento caótico. 
	
	%\coffeestainA{0.1}{0.85}{-25}{5cm}{1.3cm}
	
	
	
	\section*{Objetivos}

	\begin{itemize}
		\item Replicar integración númerica de las ecuaciones de movimiento del patinador para el caso integrable ($A_1=0$) y reproducir las gráficas de velocidad lineal y angular, y la trayectoria de la cuchilla sobre el hielo.
		\item Replicar las simulaciones de la integración númerica para el caso no integrable ($A_1=10 cm$) y reproducir la gráfica del espacio de velocidades.
		\item Replicar el cálculo de la distancia ($\delta(t)$) entre trayectorias con condiciones iniciales cercanas y ajustar la pendiente de $\log \delta(t)$ para evidenciar su divergencia exponencial.
		\item Reproducir la estimación del exponente de Lyapunov $\lambda$ para distintos valores de $A_1$.
	\end{itemize}
	
	

	\section*{Motivación}
La motivación principal de este proyecto viene de la gran fascinación que tengo por el patinaje artístico, ya que me interesa el estudio de la mecánica del movimiento del deporte y su conexión con el caos. Lo que me llamó la atención de este paper fue el manejo de un sistema complejo que es integrable en un caso particular, y que se vuelve caótico al perturbar ese caso. Al replicar los resultados, espero entender mejor lo que hay detrás del concepto de caos y cómo este se relaciona con el movimiento real del patinador. 
	\section*{Conexión con la Mecánica Clásica}
	Este proyecto conecta con la Mecánica Clásica a través de los siguientes conceptos.
	Primero, se revisan las restricciones holónomas y no holónomas: el contacto
	continuo con el hielo y el ángulo constante de la cuchilla son holónomas, mientras
	que la condición de que la cuchilla no se desliza lateralmente es no holónoma, lo
	que hace el problema más complejo que uno tratado con el formalismo de Lagrange
	estándar.
	Segundo, el patinador se modela como un cuerpo rígido, por lo que sus momentos de
	inercia son constantes.
	Tercero, las ecuaciones de movimiento describen el sistema con las velocidades $\Omega_1$, $\Omega_2$ y $v$,
	que se integrarán numéricamente.
	Cuarto, cuando la proyección del centro de masa sobre la dirección de la cuchilla
	coincide con el punto de contacto ($A_1 = 0$), el torque debido a la fuerza de
	restricción desaparece y existen constantes de movimiento adicionales, $J_1$ y
	$J_2$, que hacen el sistema integrable.
	Cuando $A_1 \neq 0$, estas cantidades dejan de conservarse, solo se conserva la
	energía $E$, y el sistema presenta comportamiento caótico \cite{gzenda2020}.
	
\section*{Metodología y procedimiento}

Planeo llevar a cabo mi proyecto estudiando las ecuaciones de movimiento del
modelo del patinador. También implementaré un programa para la integración
numérica del caso integrable, con el fin de comparar los resultados con las
gráficas de la fuente principal~\cite{gzenda2020}. Para comprobar que la
simulación es correcta, verificaré que la energía $E$ y las cantidades $J_1$ y
$J_2$ se conserven dentro del error numérico cuando $A_1 = 0$, y contrastaré
las curvas obtenidas con las del artículo. Repetiré el procedimiento para el
caso no integrable, con $A_1 \neq 0$, y analizaré las cantidades conservadas en
ese caso: se espera que $E$ se siga conservando y que $J_1$ y $J_2$ dejen de
hacerlo. Además, calcularé la divergencia entre soluciones cercanas para
estimar el exponente de Lyapunov $\lambda$. Finalmente, organizaré todo en un
repositorio de \textit{GitHub} para que las reproducciones queden al alcance
de las personas interesadas.



\section*{Material de entrega}
	


El producto final es un repositorio de \textit{GitHub} que contendrá el código en
Python desarrollado para integrar las ecuaciones de movimiento del patinador y
reproducir las gráficas de la fuente principal~\cite{gzenda2020}, junto con las
figuras obtenidas, y un archivo README con
instrucciones para ejecutar el código. Estos resultados se presentarán en el
simposio en forma de presentación comparando las gráficas obtenidas con las del artículo \cite{rocha2026symposium}.
	
	
	
	
	
	
	
	
	
	
	
	
	
	\bibliographystyle{plain}
	\bibliography{referencias}
\end{document}
